% Lösungen Einheit 5 von 5 – Stationär und langfristig (zusammenbau v0.1)
\einheitenkopf{Einheit 5 von 5 · Stationär und langfristig}

\erg{17}{a) drei Schritte vorwärts \quad b) $(600 | 300)$ aus $-0{,}2x + 0{,}4y = 0$ und $x + y = 900$ \quad c) $(0 | 500)$: B gibt nichts ab, auf lange Sicht landet alles in B \quad d) $p = 0{,}3$ (aus $0{,}8 \cdot 300 + 200p = 300$) \quad e) je 80: von A nach B $0{,}1 \cdot 800 = 80$, von B nach A $0{,}2 \cdot 400 = 80$ – Gleichgewicht, die Verteilung bleibt \quad f) $M^2 = ((0 | 2 | 0), (0 | 0 | 2), (0{,}25 | 0 | 0))$ und $M^3 = E$: nach drei Schritten ist jede Verteilung wieder die alte \quad g) Faktor $= 1{,}5$: $M \cdot v = (9 | 4{,}5) = 1{,}5 \cdot (6 | 3)$ \quad h) Für $b = 0$ springen Rot und Gelb sofort auf Blau; für b zwischen 0 und $\frac{1}{3}$ wechseln Rot und Gelb eine Zeit lang, aber je Sekunde fließt der Anteil $1 - 3b$ nach Blau, langfristig leuchtet nur Blau (Grenzmatrix); für $b = \frac{1}{3}$ fließt nichts nach Blau, Rot und Gelb wechseln endlos, Blau bleibt Blau}
\erg{18}{a) Nein: das Modell liefert das $1{,}2^{50} \approx 9100$-Fache; auf einer begrenzten Insel fehlen dafür Platz und Nahrung, die Übergänge ändern sich – das Modell taugt nur kurzfristig}
\erg{19}{a) Fehler: nur bis auf Vielfache bestimmt, die Gesamtzahl fehlt. Richtig: $(800 | 200)$ aus $x = 4y$ und $x + y = 1000$ \quad b) Bei Spaltensumme eins sagen beide Gleichungen dasselbe (sie legen nur das Verhältnis fest); erst die Gesamtzahl bestimmt, welches Vielfache gemeint ist \quad c) ja: langfristig A $= 1200$, B $= 800$ Kunden (aus $0{,}1x = 0{,}15y$ und $x + y = 2000$)}
